\documentclass[11pt, a4paper, twoside]{article}

% --- Packages ---
\usepackage[utf8]{inputenc}
\usepackage{amsmath, amssymb, amsthm}
\usepackage{geometry}
\usepackage{tcolorbox}
\usepackage{tikz-cd}
\usepackage{fancyhdr}

% --- Page Layout & Styling ---
\geometry{
	top=1in, 
	bottom=1in, 
	left=1.25in, 
	right=1.25in
}

\pagestyle{fancy}
\fancyhf{}
\fancyhead[LE,RO]{\thepage}
\fancyhead[LO]{\small \textit{Practice Problems: Riemann Surfaces \& Algebraic Curves}}
\fancyhead[RE]{\small \textit{Cohomology \& Code Constructions}}
\renewcommand{\headrulewidth}{0.4pt}

%% Custom box for problem statements
%\newtcolorbox{problembox}[1][]{
%	colback=blue!5!white,
%	colframe=blue!75!black,
%	fonttitle=\bfseries,
%	title=#1,
%	arc=3pt,
%	boxrule=1pt,
%	left=10pt, right=10pt, top=10pt, bottom=10pt
%}

% Custom box for problem statements (Blue)
\newtcolorbox{problembox}[1][]{
	colback=blue!5!white,
	colframe=blue!75!black,
	fonttitle=\bfseries,
	title=#1,
	arc=3pt,
	boxrule=1pt,
	left=10pt, right=10pt, top=10pt, bottom=10pt
}

% Custom box for solutions (Green)
\newtcolorbox{solutionbox}[1][]{
	colback=green!5!white,
	colframe=green!60!black,
	fonttitle=\bfseries,
	title=#1,
	arc=3pt,
	boxrule=1pt,
	left=10pt, right=10pt, top=10pt, bottom=10pt
}

% --- Macros ---
\newcommand{\C}{\mathbb{C}}
\newcommand{\F}{\mathbb{F}}
\newcommand{\X}{\mathcal{X}}
\newcommand{\cO}{\mathcal{O}}
\newcommand{\dx}{\,dx}

\begin{document}
	
	% ==========================================
	% COVER PAGE
	% ==========================================
	\begin{titlepage}
		\centering
		\vspace*{2in}
		
		{\Huge \bfseries Riemann Surface \& \\ Algebraic Curve Theory \par}
		\vspace{0.5in}
		{\Large \textit{A Practice Notebook for Understanding Elliptic and Hyperelliptic Codes} \par}
		
		\vspace{1.5in}
		
		{\large \textbf{Topics Covered:} \par}
		\vspace{0.2in}
		\begin{itemize}
			\item Divisors and the Riemann-Roch Theorem
			\item Sheaf Cohomology and the Mittag-Leffler Problem
			\item Meromorphic Differentials and Serre Duality
			\item Construction of the Parity-Check Matrix via Residues
		\end{itemize}
		
		\vfill
		
		{\large \today \par}
	\end{titlepage}
	
	\newpage
	
	% ==========================================
	% PROBLEM 1
	% ==========================================
	\section*{Problem 1: Differentials and the Parity-Check Matrix}
	
	\begin{problembox}[Constructing $\Omega^1(E-D)$ on a Hyperelliptic Curve]
		Let $\X$ be a hyperelliptic curve of genus $g=2$ defined over $\C$ by the affine equation:
		$$ y^2 = f(x) = x^5 - x $$
		Let $P_\infty$ denote the single point at infinity. We define our bounding divisor as $D = 3P_\infty$. 
		
		We select $n=4$ distinct finite evaluation points $E = P_1 + P_2 + P_3 + P_4$, where $P_i = (x_i, y_i)$, such that $y_i \neq 0$.
		
		\vspace{0.5em}
		\textbf{Tasks:}
		\begin{enumerate}
			\item \textbf{Holomorphic Basis:} Write down the explicit basis for the space of global holomorphic 1-forms, $\Omega^1(\X)$. Verify that the dimension agrees with the genus.
			\item \textbf{Dimension of the Syndrome Space:} Using the Riemann-Roch theorem, calculate the dimension of the space of parity-check differentials, $\dim \Omega^1(E-D)$.
			\item \textbf{Constructing the Differentials:} Find an explicit basis for the meromorphic differentials $\Omega^1(E-D)$. \textit{(Hint: Start with the holomorphic differentials and multiply by rational functions in $x$ and $y$ to introduce simple poles exactly at the points $P_i$, while ensuring the pole order at $P_\infty$ does not violate the condition imposed by $-D$.)}
			\item \textbf{The Parity-Check Matrix:} Let $\omega_1, \dots, \omega_{n-k}$ be your basis from Part 3. Express the $(n-k) \times n$ parity-check matrix $H$ explicitly in terms of the residues $\text{Res}_{P_j}(\omega_i)$. Evaluate these residues using the local coordinate $t = x - x_j$.
		\end{enumerate}
	\end{problembox}
	
	\vspace{1.5em}
	\noindent\textbf{Space for Calculations:}
	
	% The rest of the page is intentionally left blank for the student.
	\newpage
	
	% ==========================================
	% SOLUTION TO PROBLEM 1
	% ==========================================
	\section*{Solution Key: Problem 1}
	
	\begin{solutionbox}[Step-by-Step Solution: Hyperelliptic Differentials]
		\textbf{Part 1: Holomorphic Basis}\\
		For a hyperelliptic curve given by $y^2 = f(x)$ of degree $2g+1$, the basis for the space of global holomorphic 1-forms $\Omega^1(\X)$ is given by:
		$$ \omega_i = \frac{x^i \,dx}{y} \quad \text{for } i = 0, 1, \dots, g-1 $$
		Since our curve has genus $g=2$, the basis consists of exactly two differentials:
		$$ \omega_0 = \frac{dx}{y}, \quad \omega_1 = \frac{x \,dx}{y} $$
		This confirms $\dim \Omega^1(\X) = g = 2$.
		
		\vspace{1em}
		\textbf{Part 2: Dimension of the Syndrome Space}\\
		By the Riemann-Roch Theorem, the dimension of the code is $k = \deg(D) + 1 - g = 3 + 1 - 2 = 2$. 
		Because the ambient space has dimension $n=4$, the syndrome space has dimension $n - k = 4 - 2 = 2$. 
		Thus, we need to find exactly $2$ linearly independent meromorphic differentials for $\Omega^1(E-D)$.
		
		\vspace{1em}
		\textbf{Part 3: Constructing the Differentials}\\
		We need differentials with simple poles at the $n=4$ points of $E = \sum P_i$. Let $V(x) = \prod_{i=1}^4 (x - x_i)$ be the vanishing polynomial of the $x$-coordinates of our evaluation points. Dividing our holomorphic basis by $V(x)$ introduces the required poles. 
		Let:
		$$ \eta_0 = \frac{dx}{V(x)y} \quad \text{and} \quad \eta_1 = \frac{x \,dx}{V(x)y} $$
		These forms have simple poles at each $P_i = (x_i, y_i)$. (A rigorous check at $P_\infty$ using the local parameter $t = x^{-1/2}$ verifies they do not violate the $-D = -3P_\infty$ pole constraint).
		
		\vspace{1em}
		\textbf{Part 4: The Parity-Check Matrix}\\
		The parity-check matrix $H$ is constructed by evaluating the residues of our basis $\eta_0, \eta_1$ at each point $P_j$. 
		Using the local coordinate $t = x - x_j$, we find $dt = dx$. The residue of $\frac{x^i \,dx}{V(x)y}$ at $P_j$ is simply the evaluation of the non-singular part at $x = x_j$:
		$$ \text{Res}_{P_j}(\eta_i) = \frac{x_j^i}{V'(x_j) y_j} $$
		Therefore, the $2 \times 4$ parity-check matrix $H$ is:
		$$ 
		H = 
		\begin{pmatrix}
			\frac{1}{V'(x_1) y_1} & \frac{1}{V'(x_2) y_2} & \frac{1}{V'(x_3) y_3} & \frac{1}{V'(x_4) y_4} \\[8pt]
			\frac{x_1}{V'(x_1) y_1} & \frac{x_2}{V'(x_2) y_2} & \frac{x_3}{V'(x_3) y_3} & \frac{x_4}{V'(x_4) y_4}
		\end{pmatrix}
		$$
	\end{solutionbox}
	
	\newpage
	
	% ==========================================
	% PROBLEM 2
	% ==========================================
	\section*{Problem 2: Elliptic Curves and Intersections}
	
%	\begin{problembox}[Divisors of Functions on a Torus ($g=1$)]
		Let $\X$ be an elliptic curve defined by the Weierstrass equation:
		$$ y^2 = x^3 + ax + b $$
		This curve has genus $g=1$ and is topologically a Torus. The single point at infinity is denoted $\cO$ (which serves as the identity element for the curve's group law). 
		
		Let $D = 3\cO$.
		
		\vspace{0.5em}
		\textbf{Tasks:}
		\begin{enumerate}
			\item \textbf{Riemann-Roch Dimension:} Calculate the dimension $k$ of the function space $\mathcal{L}(3\cO)$.
			\item \textbf{Basis of Functions:} By examining the pole orders of $x$ and $y$ at the point $\cO$, prove that the set $\{1, x, y\}$ forms a basis for $\mathcal{L}(3\cO)$.
			\item \textbf{Line Intersections:} Let $L(x,y) = \alpha x + \beta y + \gamma = 0$ be a straight line in the affine plane. Explain geometrically why the principal divisor of the function $f = \alpha x + \beta y + \gamma$ on the curve $\X$ takes the form:
			$$ (f) = P + Q + R - 3\cO $$
			where $P, Q, R$ are the three points of intersection between the line and the elliptic curve.
			\item \textbf{The Group Law:} Use the result from Task 3 to briefly explain the geometric group law of the elliptic curve (i.e., why $P + Q + R = \cO$ in the group operation).
		\end{enumerate}
%	\end{problembox}
	
	\vspace{1.5em}
	\noindent\textbf{Space for Calculations:}
	
\end{document}